% Generated draft: requires mathematical and source review.
\documentclass[11pt]{article}
\usepackage[margin=25mm]{geometry}
\usepackage{fontspec}
\setmainfont{Latin Modern Roman}
\usepackage{amsmath,amssymb,mathtools,unicode-math}
\setmathfont{Latin Modern Math}
\usepackage{enumitem,needspace,xurl,hyperref}
\setlength{\parindent}{0pt}
\setlength{\parskip}{5pt}
\setlength{\emergencystretch}{3em}
\newcommand{\N}{\mathbb{N}}
\newcommand{\Z}{\mathbb{Z}}
\newcommand{\Q}{\mathbb{Q}}
\newcommand{\R}{\mathbb{R}}
\newcommand{\C}{\mathbb{C}}
\newcommand{\F}{\mathbb{F}}
\newcommand{\A}{\mathbb{A}}
\newcommand{\PP}{\mathbb{P}}
\newcommand{\E}{\mathbb{E}}
\newcommand{\eps}{\varepsilon}
\newcommand{\dd}{\,\mathrm{d}}
\begin{document}
\section*{UnsolvedMath ID: 6700081}
\textbf{Source problem code:} catalogue entry
\subsection{Definitions and mathematical statement}
Under the usual fundamental-class interpretation, a rationally essential \(n\)-dimensional pseudomanifold \(X\) has a rational fundamental class \([X]\) and a classifying map \(c:X\to B\pi_1(X)\) such that
\[
 c_*[X]\ne0\quad\text{in }H_n(B\pi_1(X);\mathbb Q).
\]
The source asks for such a pseudomanifold, or manifold, with an arbitrary metric, and denotes its universal cover by \(\widetilde X\). The exact compactness, orientation and pseudomanifold conventions underlying this definition need source confirmation.

With the intended lifted metric and ``Hilbert volume'', the Non-Riemannian Guth--Geroch conjecture is
\[
 \forall R>0\ \exists\widetilde x_R\in\widetilde X:\quad
 \operatorname{Hilbert\text{-}volume}
 \bigl(B_{\widetilde X}(\widetilde x_R,R)\bigr)
 \ge\operatorname{vol}(B_{\mathbb R^n}(0,R))
 =\omega_nR^n,
\]
where \(\omega_n\) is the Euclidean unit-ball volume. The center may depend on \(R\); the statement does not require a single center working at all radii. The lower bound is non-strict and applies at every positive radius. No scalar-curvature hypothesis is inserted from the surrounding program.

The excerpt does not define Hilbert volume or its normalization, the admissible arbitrary metrics, their lift to the universal cover, or the open/closed ball convention. In particular a general compatible metric is not automatically a specified length metric with a canonical lifted length construction. Hilbert volume must not be replaced by Hausdorff or another Finsler/metric volume without the original definition.

The summary describes this as an extension of Guth--Geroch motivated by Burago--Ivanov arguments beyond the Riemannian setting and labels it open, but supplies no exact theorem or bibliographic references. These are source attributions, not a verified proof for the stated arbitrary-metric class or an independently checked current status.
\subsection{Short English statement}
Must the universal cover of every rationally essential metric pseudomanifold contain, at each radius, some ball whose Hilbert volume is at least the Euclidean ball volume? The center may vary with radius, and the specialized volume and metric-lifting conventions are missing.
\subsection{Sources}
\begin{itemize}
\item Local registry: \url{data/unsolvedmath/problems.json}, permanent UnsolvedMath ID \(6700081\), source code catalogue entry; supplied statement, background, and research summary, updated 2026-08-06.
\item Misha Gromov, \emph{101 Questions, Problems and Conjectures around Scalar Curvature} (2017), printed tag [?83], occurrence 84, PDF page 83, ``Non-Riemannian Guth--Geroch'': \url{https://www.ihes.fr/~gromov/wp-content/uploads/2018/08/101-problemsOct1-2017.pdf}. Supplied normalized PDF-text extraction accessed 2026-07-29; described as incomplete.
\item Guth--Geroch and Burago--Ivanov are named in the statement/summary without separate original-paper references or exact extension hypotheses.
\end{itemize}
\end{document}
